
\subsection{Strengths}

\begin{enumerate}
\item \textbf{Automation}: Eliminates the need for manual atom-by-atom bond specification, democratizing AF3 glycan modeling for non-experts.

\item \textbf{Accuracy}: Achieves $>$98\% stereochemistry accuracy, approaching the theoretical maximum of manual specification.

\item \textbf{Speed}: Sub-second processing time enables high-throughput applications and database-scale predictions.

\item \textbf{Extensibility}: The modular architecture allows easy addition of new monosaccharide CCD mappings as needed.

\item \textbf{Open source}: Freely available implementation facilitates community contribution and validation.
\end{enumerate}

\subsection{Limitations}

\begin{enumerate}
\item \textbf{CCD coverage}: While 28+ monosaccharides are supported, some rare or modified sugars require custom CCD templates not currently available in the PDB database. Specifically, GlcN (glucosamine) and GalN (galactosamine) lack standard CCD entries.

\item \textbf{Input format dependency}: The pipeline relies on correct IUPAC/WURCS notation; errors in input strings propagate through the system.

\item \textbf{Validation scope}: Our benchmark focuses on common mammalian glycans; bacterial and plant glycans with unusual linkages may require additional testing.

\item \textbf{AF3 dependency}: The tool is designed specifically for AF3 input format and may require modification for other structure prediction tools.

\item \textbf{No structural validation}: The pipeline generates input specifications but does not validate the final AF3 predictions against experimental structures.
\end{enumerate}

\section{Conclusions}

GlycoSMILES2BAP provides an automated solution to the stereochemistry preservation problem in AlphaFold 3 glycan modeling. By converting standard glycan notations (IUPAC-condensed, WURCS) to AF3-compatible CCD+BAP format, the pipeline achieves $>$98\% stereochemistry accuracy while reducing processing time from 30--60 minutes per structure to under 1 second.

Our benchmark validation on 50 diverse glycan structures demonstrates that automated BAP generation is both feasible and reliable. The tool successfully handles common mammalian glycans, sialic acids, and rare monosaccharides, with clear pathways for extending support to additional structures.

The glycobiology community can now leverage AF3's glycan modeling capabilities without the prohibitive barrier of manual BAP specification. This advance opens new possibilities for glycoprotein engineering, vaccine design, and systematic structural glycomics.

\noindent\textbf{Availability}: Open-source implementation available at \url{https://github.com/xiaqiang/glycosmiles2bap}

\noindent\textbf{Future directions}:
\begin{itemize}
\item Extend CCD mapping coverage for rare/modified monosaccharides
\item Integrate with Privateer for automated stereochemistry validation
\item Develop web interface for non-technical users
\item Explore compatibility with other structure prediction tools
\end{itemize}

\section*{Acknowledgments}

The authors thank the GlyTouCan consortium for maintaining the glycan structure repository, and the developers of GlyLES and GlycanFormatConverter for their open-source parsing tools. We acknowledge Huang et al.\ (2025) for identifying the stereochemistry problem in AF3 glycan modeling, which motivated this work.

\section*{Author Contributions}

\noindent\textbf{Qiang Xia}: Conceptualization, Methodology, Software development, Validation, Writing -- original draft, Writing -- review \& editing.

\section*{Conflict of Interest}

The authors declare no conflict of interest.

\section*{Data Availability}

The benchmark dataset and validation results are available in the supplementary materials.

